
\documentclass{article}
\usepackage{colortbl}
\usepackage{makecell}
\usepackage{multirow}
\usepackage{supertabular}

\begin{document}

\newcounter{utterance}

\centering \large Interaction Transcript for game `codenames', experiment `risk\_low', episode 2 with deepseek{-}v4{-}pro.
\vspace{24pt}

{ \footnotesize  \setcounter{utterance}{1}
\setlength{\tabcolsep}{0pt}
\begin{supertabular}{c@{$\;$}|p{.15\linewidth}@{}p{.15\linewidth}p{.15\linewidth}p{.15\linewidth}p{.15\linewidth}p{.15\linewidth}}
    \# & $\;$A & \multicolumn{4}{c}{Game Master} & $\;\:$B\\
    \hline

    \theutterance \stepcounter{utterance}  
    & & & \multicolumn{2}{p{0.3\linewidth}}{
        \cellcolor[rgb]{0.9,0.9,0.9}{
            \makecell[{{p{\linewidth}}}]{
                \texttt{\tiny{[GM$|$GM]}}
                {'hidden': {'team': ['قمر', 'ليمون', 'سيارة', 'فرنسا', 'ممرضة', 'وقت', 'كاتشب', 'فرشاة', 'كنغر'], 'innocent': ['بطارية', 'أفريقيا', 'توت', 'الأمازون', 'النينجا', 'قرصان', 'صدر', 'صوت'], 'opponent': ['الجهاز اللوحي', 'نار', 'ماء', 'زلة', 'قوس', 'موت', 'تفاحة', 'مجهر'], 'assassin': []}, 'revealed': {'team': {'team': [], 'innocent': [], 'opponent': [], 'assassin': []}, 'opponent': {'team': [], 'innocent': [], 'opponent': [], 'assassin': []}}}
            }
        }
    }
    & & \\ \\

    \theutterance \stepcounter{utterance}  
    & & \multicolumn{4}{p{0.6\linewidth}}{
        \cellcolor[rgb]{0.9,0.9,0.9}{
            \makecell[{{p{\linewidth}}}]{
                \texttt{\tiny{[P1$\langle$GM]}}
                \texttt{قدّم تلميحًا واحدًا مكوّنًا من كلمة واحدة يرتبط ببعض كلمات فريقك من القائمة أدناه. يمكنك اختيار عدد الكلمات التي يرتبط بها تلميحك وأيّها.} \\
\texttt{الهدف هو أن يخمّن شريكك جميع كلمات فريقك بناءً على التلميحات التي تقدّمها له، قبل أن يكشف الفريق الآخر جميع كلماته. إذا خَمَّن شريكك كلمة 'الخصم' أو كلمة 'المشتّت'، تنتهي جولة التخمين الخاصة به ويحصل الفريق المنافس على دور. إذا خَمَّن شريكك كلمة 'القاتل'، تنتهي اللعبة ويخسر فريقك فورًا.} \\
\texttt{كلما ارتبط تلميحك بعدد أكبر من الكلمات، كان بإمكانك الفوز أسرع، لكن قد تصبح الروابط أصعب على شريكك في التخمين، لذا اختر كلمة التلميح بحكمة. يجب أن تكون كلمة التلميح مرتبطة دلاليًا بالكلمات المستهدفة، ولا يمكن أن تكون واحدة من الكلمات في القوائم أو أن تحتوي على أجزاء منها.} \\
\\ 
\texttt{قدّم دائمًا تلميحك المكوّن من كلمة واحدة وقائمتك المفصولة بفواصل من الكلمات المستهدفة المرتبطة بالتنسيق التالي واجعل إجاباتك قصيرة قدر الإمكان، ولا تُدرج أي نص آخر غير المطلوب في هذا النموذج:} \\
\\ 
\texttt{تلميح: <كلمة>} \\
\texttt{أهداف: <أهداف>} \\
\\ 
\texttt{كلمات فريقك هي: قمر, ليمون, سيارة, فرنسا, ممرضة, وقت, كاتشب, فرشاة, كنغر.} \\
\texttt{كلمات فريق خصمك هي: الجهاز اللوحي, نار, ماء, زلة, قوس, موت, تفاحة, مجهر.} \\
\texttt{كلمات التشتيت هي: بطارية, أفريقيا, توت, الأمازون, النينجا, قرصان, صدر, صوت.} \\
\texttt{كلمات القتل هي: .} \\
            }
        }
    }
    & \\ \\

    \theutterance \stepcounter{utterance}  
    & \multicolumn{4}{p{0.6\linewidth}}{
        \cellcolor[rgb]{0.9,0.9,0.9}{
            \makecell[{{p{\linewidth}}}]{
                \texttt{\tiny{[P1$\rangle$GM]}}
                \texttt{تلميح: حامض} \\
\texttt{أهداف: ليمون، كاتشب} \\
            }
        }
    }
    & & \\ \\

    \theutterance \stepcounter{utterance}  
    & & & \multicolumn{2}{p{0.3\linewidth}}{
        \cellcolor[rgb]{0.9,0.9,0.9}{
            \makecell[{{p{\linewidth}}}]{
                \texttt{\tiny{[GM$|$GM]}}
                {'player': 'Cluegiver', 'type': 'target is hallucination', 'utterance': 'تلميح: حامض\nأهداف: ليمون، كاتشب', 'message': "Targeted word 'ليمون، كاتشب' was not listed, you can only target words provided in the lists.", 'target': 'ليمون، كاتشب', 'previous_targets': [], 'remaining_words': ['ليمون', 'ماء', 'ممرضة', 'بطارية', 'الأمازون', 'صدر', 'النينجا', 'موت', 'تفاحة', 'نار', 'كنغر', 'الجهاز اللوحي', 'فرنسا', 'قمر', 'صوت', 'مجهر', 'أفريقيا', 'سيارة', 'زلة', 'توت', 'قرصان', 'كاتشب', 'وقت', 'فرشاة', 'قوس']}
            }
        }
    }
    & & \\ \\

\end{supertabular}
}

\end{document}
